
Section~\ref{sec:results-hardware} showed that the first call to
AlphaFold2's forward pass at a new input shape is far more expensive
than steady state, and that part of that cost is an artifact of
capturing the call inside a profiler trace rather than compilation
itself (Section~\ref{sec:rh-coldstart}). This section looks inside a single
first call to identify, at the function level, where the remaining
cost goes.

\subsection{Tracing a first call}
\label{sec:prof-method}

We captured one first-\texttt{predict} call for the single-query TPU
configuration (\texttt{model\_3}, 118 residues, \texttt{num\_recycle=0})
inside \texttt{jax.profiler.trace}, copied the resulting trace off the
pod before it was reclaimed, and inspected it in Chrome's trace
viewer. TPU pods on this cluster are torn down within roughly a
minute of job completion, so capturing the trace file required an
explicit copy step while the pod was still alive.

\subsection{The \texttt{cache\_miss} finding}
\label{sec:prof-cachemiss}

The traced call stack descends from AlphaFold2's outer
\texttt{apply\_fn} wrapper, which itself accounts for essentially none
of the traced time, into a function named \texttt{cache\_miss}, defined
in JAX's own \texttt{pjit.py}. Its purpose is to trace and compile a
computation the first time JAX encounters a given input shape. Of the
16.56\,s spanned by the traced \texttt{apply\_fn} call, 12.55\,s (about
76\%) is self time inside \texttt{cache\_miss} itself, not in any
function it calls. These trace figures come from the analysis written
up when the trace was taken; the raw profiler trace is not in the
repository, so they are reported figures that cannot be re-derived
from a retained artifact (Section~\ref{sec:limitations}).
Below \texttt{cache\_miss}, the call stack descends through
\texttt{\_infer\_params}, \texttt{\_trace\_for\_jit}, and
\texttt{trace\_to\_jaxpr} into Haiku's own \texttt{apply\_fn}, which
lowers AlphaFold2's Evoformer computation into XLA's intermediate
representation. Every frame in this path executes on the host. The
TPU device track in the same trace shows only a handful of brief
marks near the start of the window, and is otherwise idle for nearly
the full 16.56\,s while the host finishes tracing and compiling.

The 76\% figure is a share of the traced span, not of the
first-\texttt{predict} call as a whole: the trace covers 16.56\,s,
while the first-\texttt{predict} timings reported in
Section~\ref{sec:results-hardware} run 27.43--28.80\,s
in the runs that measure it directly. The two are separate captures
of related but not identical work, and the trace does not by itself
explain the gap between them. We report the 76\% figure only as a
property of the traced \texttt{apply\_fn} span: three quarters of that
span is JAX's own tracing-and-compilation machinery, concentrated in a
single named function, rather than TPU execution, data movement or
AlphaFold2's own arithmetic.

\begin{figure}[H]
  \centering
  \includegraphics[width=\linewidth]{coldstart_regions.pdf}
  \caption{The cold-start regime. (a) The three timed regions of the
  August baselines on each backend, on a log scale. On CPU and GPU the
  lighter top segment of the first-\texttt{predict} bar is the profiler
  teardown that ran inside the timer (36.00\,s and 42.02\,s,
  Section~\ref{sec:rh-coldstart}); the TPU bar cannot be corrected
  because no TPU run log survives. (b) On TPU, a persistent compilation
  cache shortens \texttt{init\_params} 6.81$\times$ and the first
  \texttt{predict} call 1.90$\times$ in a fresh process; the steady state
  is unchanged.}
  \label{fig:coldstart}
\end{figure}

\subsection{Compilation cache}
\label{sec:prof-cache}

A persistent JAX compilation cache stores compiled XLA binaries across
process restarts, so a second process at the same input shape can skip
the compile step that produced them. Comparing a cold process against
one with a warm cache, on the same single-query configuration
(Figure~\ref{fig:coldstart}b),
\texttt{init\_params} drops from 37.68\,s to 5.53\,s, a 6.81$\times$
speedup, and the first \texttt{predict} call drops from 28.80\,s to
15.19\,s, a 1.90$\times$ speedup.
These are two different compilations and the two speedups should not
be compared to each other or to the 12.55\,s \texttt{cache\_miss}
figure above: the cache experiment and the trace are separate runs, and
a warm cache still requires a fresh process to re-trace the Python
program down to a jaxpr before it can consult the cache, which
plausibly accounts for part of the 15.19\,s that remains on the warm
first call. The comparison does establish that most of AlphaFold2's
start-up cost on TPU is avoidable across process restarts at a fixed
input shape, provided that cost is taken as \texttt{init\_params} and
the first \texttt{predict} together: the two fall from 66.48\,s to
20.72\,s, a 69\% reduction. The first \texttt{predict} alone falls only
47\%, so more than half of it survives the cache, and
\texttt{init\_params} benefits far more from caching than the forward
pass itself does.

\subsection{Relation to prior work}
\label{sec:prof-priorwork}

Compilation-dominated cold starts are not unique to this hardware or
this study. ParaFold identifies JAX recompilation as a major cost of
running AlphaFold2 at scale on GPU, and mitigates it within a process by
sorting input sequences by length so that fewer distinct shapes
trigger a fresh trace-and-compile cycle \citep{zhong2022parafold}. Our
finding is the same underlying cost, reported at the function level by
our retained trace analysis, on a different accelerator. We take it as consistent evidence that
the cost comes from AlphaFold2's JAX/XLA compilation model, whatever
the backend, and do not claim it as an independent discovery.
